% Fall 2 (Brandmeldeanlage): Profil, Annahmen und Prozess, Entwurf 22.11.2026
\paragraph{Profil}
\begin{small}\begin{longtable}{@{}>{\raggedright\arraybackslash}p{5.0cm}>{\raggedright\arraybackslash}p{9.5cm}@{}}
\caption[Organisationsprofil für Fall~2]{Organisationsprofil für Fall~2: Brandmeldeanlage. Quelle: eigene Darstellung (Entwurf).}\label{tab:fall2_profil}\\
\toprule \textbf{Merkmal} & \textbf{Angabe} \\ \midrule \endfirsthead
\toprule \textbf{Merkmal} & \textbf{Angabe} \\ \midrule \endhead
\bottomrule \endfoot
Name & Eventhalle Süd GmbH (fiktiv) \\
Sitz & München, Bayern, Deutschland \\
Rechtsform & Gesellschaft mit beschränkter Haftung \\
Sektor & privat \\
Branche & Veranstaltungen \\
Beschäftigte & 80 \\
Arbeitnehmervertretung & Betriebsrat \\
Gebäude & neue Veranstaltungshalle für bis zu 3.000 Personen \\
KI-Einsatz & gekauftes KI-System, das Brandmeldungen bewertet \\
Dienstleister außerhalb der EU & nein \\
\end{longtable}\end{small}
\paragraph{Annahmen}
\begin{description}[leftmargin=1.0cm,style=sameline]
  \item[A1] Die GmbH ist Bauherrin und Betreiberin der Halle.
  \item[A2] Die Baugenehmigung liegt vor. Das Brandschutzkonzept verlangt eine Brandmeldeanlage mit Aufschaltung zur Feuerwehr.
  \item[A3] Das KI-System stammt von einem Hersteller in der EU und verarbeitet nur Daten der Melder, keine Videobilder.
  \item[A4] Errichtung und Instandhaltung übernimmt ein externer Fachbetrieb.
\end{description}
\paragraph{Prozess}
\begin{small}\begin{longtable}{@{}>{\raggedright\arraybackslash}p{1.0cm}>{\raggedright\arraybackslash}p{3.6cm}>{\raggedright\arraybackslash}p{10.0cm}@{}}
\caption[Prozess für Fall~2: Brandmeldeanlage]{Prozess für Fall~2: Brandmeldeanlage mit Teilprozessen, Tasks und Beschreibungen. Quelle: eigene Darstellung (Entwurf).}\label{tab:fall2_prozess}\\
\toprule \textbf{Kennung} & \textbf{Bezeichnung} & \textbf{Beschreibung} \\ \midrule \endfirsthead
\toprule \textbf{Kennung} & \textbf{Bezeichnung} & \textbf{Beschreibung} \\ \midrule \endhead
\bottomrule \endfoot
\multicolumn{3}{@{}l}{\textbf{TP1 Planen}} \\
T1.1 & Brandschutzkonzept auswerten & Die Bauleitung liest das Brandschutzkonzept und listet die Anforderungen an die Brandmeldeanlage auf. \\
T1.2 & Anlagenkonzept erstellen & Ein Fachplaner legt Meldebereiche, Meldertypen und Alarmwege fest. \\
T1.3 & Aufschaltung abstimmen & Der Fachplaner klärt mit der Feuerwehr, wie Alarme an die Leitstelle übertragen werden. \\
T1.4 & KI-System auswählen & Die Geschäftsführung wählt ein KI-System, das Meldungen auf Täuschungsalarme prüft. \\
\multicolumn{3}{@{}l}{\textbf{TP2 Beauftragen}} \\
T2.1 & Angebote einholen & Die Bauleitung holt Angebote von drei Fachbetrieben ein. \\
T2.2 & Errichter beauftragen & Die Geschäftsführung beauftragt den günstigsten Fachbetrieb mit Errichtung und Instandhaltung. \\
\multicolumn{3}{@{}l}{\textbf{TP3 Errichten}} \\
T3.1 & Melder und Zentrale montieren & Der Errichter montiert Melder, Handfeuermelder und die Brandmelderzentrale. \\
T3.2 & Leitungen verlegen & Der Errichter verlegt die Leitungen zwischen Meldern, Zentrale und Übertragungseinrichtung. \\
T3.3 & KI-System anbinden & Der Errichter verbindet das KI-System mit der Brandmelderzentrale. \\
\multicolumn{3}{@{}l}{\textbf{TP4 Abnehmen}} \\
T4.1 & Anlage prüfen lassen & Ein Prüfsachverständiger prüft die Anlage und stellt einen Prüfbericht aus. \\
T4.2 & Aufschaltung testen & Die Feuerwehr löst Probealarme aus und prüft den Eingang in der Leitstelle. \\
T4.3 & Dokumentation übergeben & Der Errichter übergibt Pläne, Prüfprotokolle und Bedienungsanleitungen an die Betreiberin. \\
\multicolumn{3}{@{}l}{\textbf{TP5 Betreiben}} \\
T5.1 & Verantwortliche benennen & Die Geschäftsführung benennt den Haustechniker als Verantwortlichen für die Anlage. \\
T5.2 & Meldung bewerten & Das KI-System bewertet jede Meldung und stuft sie als wahrscheinlichen Brand oder als Täuschungsalarm ein. \\
T5.3 & Meldung erkunden & Bei einem vermuteten Täuschungsalarm prüft ein Mitarbeiter den Melder vor Ort innerhalb von drei Minuten. \\
T5.4 & Alarm weiterleiten & Die Zentrale leitet bestätigte Alarme und nicht erkundete Meldungen an die Leitstelle der Feuerwehr weiter. \\
T5.5 & Betriebsbuch führen & Der Haustechniker trägt jede Meldung, Störung und Abschaltung in das Betriebsbuch ein. \\
\multicolumn{3}{@{}l}{\textbf{TP6 Instandhalten}} \\
T6.1 & Anlage inspizieren & Der Fachbetrieb prüft Zentrale und Melder auf sichtbare Mängel und Funktion. \\
T6.2 & Anlage warten & Der Fachbetrieb reinigt und tauscht Melder und Akkus nach Herstellerangabe. \\
T6.3 & Wiederkehrend prüfen lassen & Ein Prüfsachverständiger prüft die Anlage erneut und stellt einen Prüfbericht aus. \\
\end{longtable}\end{small}
\offen{Referenzzuordnung für Fall~2: fehlende Schritte (Leistungserklärung der Melder, Feuerwehrpläne, Brandschutzordnung, Unterweisung, KI-Kompetenz), Kontrollschritt T4.1, Falle Hochrisiko-Einstufung nach AI Act}
